% pxGNT -- optional P3CTeX Gantt module
\documentclass[a4paper]{article}
\usepackage[english]{babel}
\usepackage{pxGNT}
\begin{document}
\section*{pxGNT: daily Gantt charts}
The optional \texttt{GNT} module provides a three-command interface to
\texttt{pgfgantt}; it does not change the document orientation.
Select it when loading P3CTeX: \verb|\documentclass[GNT,PR,CAT]{P3CTeX}|.
Outside the class, use \verb|\usepackage{pxGNT}|.

\begin{verbatim}
\pxGantt{start=01/10/2026,end=31/01/2027,days-per-block=28}
\pxGanttItem{Investigation}{01/10/2026}{20/10/2026}
\pxGanttItem{Implementation}{15/10/2026}{15/12/2026}
\pxGanttEnd
\end{verbatim}

Options: \texttt{start}, \texttt{end} (required dd/mm/yyyy),
\texttt{days-per-block} (default 28), \texttt{days-per-page} (alias),
\texttt{day-width} (optional TeX length), \texttt{label-width} (manual override),
\texttt{label-max-width} (default \verb|.24\linewidth|),
\texttt{rotate=true} (optional chart-only rotation).

By default the label column is measured from actual task text and capped
relative to the available line width. Remaining width is divided among
calendar days. The date range is balanced into consecutive blocks, and
bars crossing block boundaries are clipped. P3CTeX colours are reused;
standalone users receive fallback colours.

\section*{Example}
\pxGantt{start=01/10/2026,end=25/11/2026,days-per-block=28}
\pxGanttItem{Investigació}{01/10/2026}{20/10/2026}
\pxGanttItem{Disseny}{15/10/2026}{10/11/2026}
\pxGanttItem{Implementació}{01/11/2026}{25/11/2026}
\pxGanttEnd
\end{document}
